\section{Framework Overview}

PrismMath Tutor is organized around a simple principle: the language model should remain responsible for explanation, dialogue, and pedagogical framing, while external components provide context, persistence, and exact computation. The system is not a monolithic prompt. It is a small agentic framework whose behavior is assembled from retrieval, memory, tool use, logging, and a shared chat-turn routine. This section describes the deployed framework before Section~\ref{sec:finetuning} describes the AVSD training method for the underlying reasoning model.

\subsection{Agentic Components}

The current codebase exposes two front ends over the same core logic. The command-line demo starts from \texttt{main.py}, delegates to \texttt{app.chat\_cli}, prompts for a student profile, and then loops over student questions. The web demo is served by \texttt{app.web\_ui} through FastAPI and renders \texttt{app/templates/chat.html}. Both front ends call \texttt{run\_chat\_turn} in \texttt{app.chat\_core}, which performs the same sequence: load recent memory, optionally retrieve examples, build a student-facing prompt, call the LLM client, parse any tool request, execute the tool, ask the model for a final answer after tool execution, and persist the result.

This organization matters because tutoring failures are often system failures rather than only model failures. If a student asks a question that resembles examples already present in the corpus, the tutor should have access to those examples. If the same student returns for another turn, the tutor should remember recent questions and the preferred explanation style. If the model needs to solve a large equation or verify a fraction, it should not rely only on free-form arithmetic in generated text. The agentic components are therefore designed to solve distinct failure modes.

\subsubsection{Retrieval-Augmented Generation}

The RAG component addresses the problem of mathematical grounding \citep{lewis2020rag}. A general-purpose language model can often explain algebraic procedures, but it may not remember the exact style, difficulty, or solution patterns of the target corpus. PrismMath uses a local sample of MetaMathQA-style examples as an in-context source. The expected corpus path is \texttt{data/processed/metamathqa\_30k.jsonl}. Each record must contain at least an identifier, a question, and an answer. The loader in \texttt{app.rag.dataset} validates these fields and returns a dictionary keyed by example identifier.

Index construction is handled by \texttt{app.rag.ingest}. The ingestion routine loads the local examples, encodes their questions with a local Qwen embedding model, and stores the resulting vectors in a persistent ChromaDB collection. The configured collection name is \texttt{metamathqa\_questions\_qwen3\_4b}; the index is stored under \texttt{data/chroma}. The embedding model is expected under \texttt{models/qwen3-embedding-4b/}. This local-model requirement is intentional: retrieval quality should not depend on a remote embedding endpoint changing over time, and the capstone demo can be rebuilt from declared local artifacts.

At runtime, \texttt{RAGRetriever} encodes the student's first question, queries ChromaDB for the top-$k$ nearest examples, and returns \texttt{RetrievedExample} objects containing the example id, question, answer, and similarity score. The prompt builder then serializes these examples into an in-context block. The examples are not treated as ground truth for the student's exact question. The system prompt explicitly tells the model to use retrieved MetaMathQA examples as in-context examples while answering the current question. This distinction is important: RAG supplies analogies and solution style, not a license to copy an answer for a different problem.

RAG solves three practical problems in PrismMath. First, it gives the tutor examples of mathematical exposition that are close to the student's query. Second, it reduces dependence on model memorization by making local corpus information visible in the prompt. Third, it makes the demo inspectable: the web interface displays the top retrieved examples, allowing evaluators to see whether the retrieval context was relevant. This is particularly useful in an educational setting, where a plausible but irrelevant example can mislead the explanation.

The current implementation retrieves examples only on the first turn of a browser or CLI session. Later turns use conversation memory instead. This is a conservative design choice. First-turn retrieval ensures that the tutor starts with domain-specific support when the student introduces a problem. For follow-up turns, recent interaction history is often more relevant than another independent nearest-neighbor query. The limitation is that a later turn introducing a completely new problem will not automatically trigger retrieval unless the session is restarted or the retrieval policy is extended. A future version should support intent-aware retrieval on any turn that shifts to a new problem.

\subsubsection{Student Memory}

Student memory addresses personalization and continuity. A tutor should not treat every question as the first interaction with an anonymous user. PrismMath stores profile fields such as student id, nickname, grade, weak areas, strong areas, and preferred explaining style. It also stores previous interactions, retrieved example ids, and tool calls. The memory layer is implemented in \texttt{app.memory.store} with SQLite and writes to \texttt{data/demo.sqlite3}.

The profile is created or updated at the start of a CLI or web session. In the web interface, form fields are submitted with each question, and \texttt{MemoryStore.upsert\_profile} persists the latest version. In the CLI, the program first attempts to load an existing profile and only asks for profile fields if no profile exists. This gives the demo a simple but useful form of identity continuity: the same student id can recover prior preferences without requiring a separate authentication system.

Conversation continuity is implemented through recent interaction retrieval. Before each model call, \texttt{run\_chat\_turn} asks the store for the latest interactions for that student. The prompt builder formats them as a recent conversation block. This memory is short-range rather than a full long-term learner model. It is sufficient for continuity across a tutoring session: the model can refer to what the student just asked, avoid repeating the same explanation, and adapt phrasing to the stated preferred style. It also avoids the complexity and privacy risk of storing a broad inferred psychological profile.

Memory also improves auditability. Each interaction stores the query, answer, retrieved example ids, and serialized tool calls. The web and CLI sessions additionally create text logs through \texttt{ConversationLogger}. These logs include the student profile, each question, retrieved examples, tool calls, and final answer. For a capstone demonstration, this is valuable because one can inspect a complete turn after the fact: what the student asked, what context the model saw, whether a tool was used, and what answer was returned.

The memory design intentionally separates user data from source code. The README documents deletion commands for removing a single student's profile, interactions, tool-call rows, retrieved-example rows, and logs, as well as commands for deleting all demo user details. This matters for educational systems because even a prototype can store sensitive learning information. The report therefore treats \texttt{data/demo.sqlite3} and \texttt{logs/} as generated local artifacts, not as files to be committed or shared.

There are several limitations. The current memory is mostly extractive: it stores recent raw interactions and has a table for summaries, but the active prompt path uses recent turns rather than a learned student model. It does not yet infer mastery, misconception categories, or curriculum state. It also does not implement memory compaction beyond selecting a fixed number of recent interactions. These are reasonable constraints for a prototype. They keep the system transparent and reduce the risk that hidden long-term inferences dominate the tutor's behavior.

\subsubsection{Symbolic Verification Tool}

The symbolic verification tool addresses a different failure mode: exactness. Math tutoring systems often fail not because the reasoning plan is wrong, but because a generated intermediate calculation is off by a small amount. PrismMath gives the model a single tool, \texttt{verify\_answer}, implemented with SymPy \citep{meurer2017sympy}. The tool can simplify expressions, solve equations, compare a proposed answer to a symbolic result, and support derivative and integral helper functions. The public prompt contract, however, exposes only the \texttt{verify\_answer} action.

Tool use begins in the system prompt. The model is instructed to decide internally whether a question is \emph{DIRECT} or \emph{TOOL}. Direct questions can be answered without external computation. Tool questions are those where exact computation, symbolic simplification, equation solving, derivatives, integrals, large arithmetic, or nontrivial fractions would materially improve correctness. If the model chooses tool use, it must output only a fenced block containing JSON with the action and arguments. This strict contract makes the tool call easy to parse and prevents the first response from mixing natural language with a half-formed tool request.

The parser in \texttt{app.tools.runtime} searches for a fenced \texttt{tool} block, decodes the JSON, and rejects unsupported tool names. Dispatch is deliberately narrow: only \texttt{verify\_answer} is accepted. The tool implementation sanitizes input with a restricted character pattern and converts caret notation to Python exponentiation. For equations, it constructs a SymPy equality, solves for the available symbol, and optionally verifies a provided answer. For expressions, it simplifies the expression and optionally compares the simplified result to an expected answer. After tool execution, \texttt{run\_chat\_turn} builds a follow-up prompt containing the original student question, the draft tool request, and the tool result. The model then produces a final student-facing answer.

This two-step pattern supports both correctness and pedagogy. The tool returns a compact symbolic result, but it does not write the explanation. The language model converts that result into a step-by-step answer in the student's language and preferred style. For example, for the equation
\[
987654321x - 123456789 = 555555555,
\]
the tool can verify the symbolic solution while the final answer can explain that adding $123456789$ to both sides gives $987654321x=679012344$, and reducing the fraction gives $x=75445816/109739369$. The tool protects the arithmetic; the model supplies the tutoring.

The tool is intentionally constrained. Its input contract forbids natural language in the \texttt{problem} argument and requires explicit multiplication such as \texttt{987654321*x}. This reduces parsing ambiguity and limits unsafe symbolic input. The tradeoff is that the model must learn to produce clean mathematical strings. When it fails, the tool parser or SymPy layer can reject the request, and the tool-call record captures the error. In a production tutor, this tool layer would likely be expanded with structured schemas for equation solving, arithmetic, graphing, and proof checking. For the current prototype, one explicit verifier is enough to demonstrate agentic tool use without obscuring the main system.

\subsection{Multi-context Self-Distillation}

The framework components above determine how PrismMath behaves at inference time. The remaining question is how to train or adapt the model that performs the reasoning. PrismMath adopts the view that the reasoning model should learn from multiple training-time contexts, not just from one static answer trace. A math dataset can often provide a full solution, a partial rationale, a final answer, or generated feedback. These views are not all equally useful. A full solution may be informative but can force the model toward a reference reasoning path. A final answer may avoid over-constraining the explanation but gives less token-level guidance. A partial rationale can encourage the early structure of a solution while leaving later steps flexible.

AVSD provides a principled way to combine these views. Instead of choosing one privileged context as the teacher, it evaluates a family of view-conditioned teachers on the same student-generated prefix. Signals that agree across views are treated as stable consensus. Signals that appear only in one or a few views are treated as residual support and are added only when they align with the consensus and remain proportionate. In PrismMath terms, AVSD is the training-side analogue of the runtime architecture: multiple sources of context are useful, but each must be controlled so that it helps the tutor without dominating the final behavior.
